\documentclass[parskip=full]{scrartcl}
\usepackage{geometry}
\geometry{a4paper, left=2.5cm, right=2cm, top=3cm, bottom=3cm}
\usepackage[utf8]{inputenc} % use utf8 file encoding for TeX sources
\usepackage[T1]{fontenc}    % avoid garbled Unicode text in pdf
\usepackage[german]{babel}  % german hyphenation, quotes, etc
\usepackage{hyperref}       % detailed hyperlink/pdf configuration
\hypersetup{                % ‘texdoc hyperref‘ for options
pdftitle={Echtzeit-Computergrafik in der Spieleentwicklung},%
bookmarks=true,%
}

\usepackage{xcolor}
\usepackage{graphicx}       % provides commands for including figures
\usepackage{csquotes}       % provides \enquote{} macro for "quotes"
\usepackage[nonumberlist]{glossaries}     % provides glossary commands
\usepackage{enumitem}
\usepackage{epsf,epsfig,eepic}
\usepackage[headsepline, footsepline]{scrpage2}

%kopf und fusszeilen
\pagestyle{scrheadings}
\clearscrheadfoot

%kopfzeile
\ihead{\codename}
\ohead{\headmark}
\automark{section}

%fusszeile
\ofoot{\pagemark}
\ifoot{\codename}

\makenoidxglossaries
%
% % Glossareinträge
%

\newglossaryentry{KI}{
	name={KI},
	plural={KI's},
	description={Eine Künstliche Intelligenz (KI) beschreibt hier einen automatisierten, nicht menschlich gesteuerten, Fahrer},
}
\newglossaryentry{PBR}{
	name={PBR},
	description={Physically Based Rendering},
}
\newglossaryentry{Deferred Rendering}{
	name={Deffered Rendering},
	description={beschreibt eine Technik, mit deren Hilfe die
	           Geometrieverarbeitung von der Lichtberechnung getrennt werden kann}
}
\newglossaryentry{MotionBlur}{
	name={MotionBlur},
	description={Unschärfe von bewegten Objekte in Bildern}
}

\newglossaryentry{Seed}{
	name={Seed},
	plural={Seeds},
	description={Seed (zu deutsch: Samen) beschreibt einen Ausgangswert zur Generierung einer Folge von Pseudozufallswerten},
}

\newglossaryentry{Biom}{
	name={Biom},
	plural={Biome},
	description={Ein Biom beschreibt eine bestimmte Ausprägung eines Streckenabschnittes},
}

\newglossaryentry{Mesh}{
	name={Mesh},
	plural={Meshes},
	description={Polygonnetz welches in der Computergrafik zur Darstellung von Formen im Raum dargestellt werden},
}

\newglossaryentry{Navigation Mesh}{	
	name={Navigation Mesh},
	plural={Navigation Meshes},
	description={Repräsentiert die von Spieler (und KI) befahrbare Oberfläche}
}

\newglossaryentry{Scenegraph}{
	name={Scenegraph},
	description={Baumartige Datenstruktur zur Repräsentation einer Szene}
}
\newglossaryentry{JUnit}{
	name={JUnit},
	description={JUnit ist ein Werkzeug welches Quelltext, mit vorher festgelegten Testflällen, automatisiert testet}
}
\newglossaryentry{Low-Poly Stil}{
	name={Low-Poly Stil},
	description={Low-Poly ist ein Polygonnetz in der 3D-Computergrafik, welches eine relativ kleine Anzahl von Polygonen besizt}
}
\newglossaryentry{FXAA}{
	name={Fast Approximate Anti-Aliasing (FXAA)},
	description={Reduziert sichtbares Aliasing während Schärfe beibehalten wird}
}
\newglossaryentry{HDR}{
	name={High Dynamic Range (HDR)},
	description={Ist das Rendern von virtuellen Umgebungen mit einem dynamischen Farbkanalbereich von 65.535: 1 oder höher (in der Computer-, Spiel- und Unterhaltungstechnologie)}}
%Kommandos
\newcommand{\codename}{Valaris}

%
% % Abkürzungen
%
%TODO

\title{Pflichtenheft}
\author{Sidney Hansen, Jonas Heinle, Frederik Lingg, Lukas Schölch, Artur Wesner}

\begin{document}

	\begin{titlepage}
		\centering
		{\scshape\LARGE Pflichtenheft Praxis der Softwareentwicklung\par}
		\vspace{1cm}
		{\scshape\Large Echtzeit Computergrafik in der Spieleentwicklung \par}
		\vspace{1cm}
		{\huge\bfseries \codename \par}
		\vspace{1cm}
		\includegraphics[width=.5\linewidth]{./Bilder/valaris_t_4096.png}
		\par
		{\vspace{1cm}}
		{\Large\itshape Autoren \\}
		{\Large\itshape Sidney Hansen, Jonas Heinle,\\}
		{\Large\itshape Frederik Lingg, Lukas Schölch, Artur Wesner\par}
		
		\vfill
		Projektbetreuer\par
		Alisa Jung \& Tobias Rapp
		
		\vfill
		
		% Bottom of the page
		{\large \today\par}
	\end{titlepage}

	\tableofcontents 
	
	\pagebreak
	
	\section{Einleitung}
	
	%Benutzer-Erfahrung
	\codename\ ist ein lineares futuristisches Weltraumrennspiel. Spieler können gegen andere Spieler oder virtuelle Gegner in kontrollpunktbasierten Wettrennen antreten.
	Das Rennen wird auf schwebenden Fahrzeugen bestritten. Die Strecke führt durch eine prozedural generierte Planetbasisstation, welche abwechslungsreiche Umgebungen
	für spannende Wettkämpfe bietet. Während des Rennens können Spieler Energie-Kugeln klauen um schneller als die Konkurrenz zu sein. An Kontrollpunkten erhalten Spieler
	abhängig von ihrer Position Energie-Kugeln, um die Anzahl der Energie-Kugeln der Spieler auszugleichen. Zusätzlich sind Powerups auf der Strecke verteilt, um den Spielern
	weitere taktische Möglichkeiten zur Verfügung zu stellen. \par

	%Entwicklungs Ziele
	\textcolor{gray}{[Team Grafik]}\\
	\codename\ soll als Echtzeit-Java-Anwendung realisiert werden. Die Render-Pipeline soll für andere Projekte wiederverwendbar sein. Für die Strecken soll ein modularer
	Generator entwickelt werden, sodass das Hinzufügen neuer Features für die Strecke mit minimalem Aufwand möglich ist. Weiter soll eine intuitive Menüführung für
	\codename\ erstellt werden, welche es dem Nutzer ermöglicht die Steuerung auf seine Vorlieben und die Grafik auf sein System anzupassen. \par

	\textcolor{gray}{[Team Gameplay]}\\
	\codename\ simuliert den oben beschriebenen Rennablauf. Hierzu wird ein physikalisches Modell des Rennens gebildet. In diesem Modell wird das Fahrverhalten
	des Fahrzeuges berechnet. Indem Kollisionen, unterschiedliche Streckenuntergründe und die Steuerung des Spielers die Berechnungen beeinflussen.
	Dies soll sowohl lokal (Einzelspieler) als auch über Netzwerkverbindung (Mehrspieler) erfolgen. \par

	%Programmablauf
	\textcolor{gray}{[Team Grafik]}\\
	Nach dem Start von \codename\ befindet sich der Nutzer in einem Hauptmenü, von dem aus er entweder zu den Einstellungen oder zum Spiel navigieren kann. Vor dem Start eines Rennens
	können die Spieler verschiedene Rennwagen aussuchen, währenddessen wird das Level im Hintergrund geladen. Danach startet das Rennen.
	Nach dem Rennen können die Spieler gegenseitig ihre Leistung ansehen, woraufhin sie zurück ins Hauptmenü gelangen. \par

	\textcolor{gray}{[Team Gameplay]}\\
	Das Rennen startet an einer Startlinie mit einem visuellen Signal. Während des Rennens kann das Fahrzeug mit einem Kontroller 
	oder einer Tastatur gesteuert werden. Jeder Spieler startet mit X Energie-Kugeln. Die Energiekugeln erhöhen die 
	Maximalgeschwindigkeit des Fahrzeuges. Diese Energiekugeln kann man von Gegenspielern klauen. 
	Zudem erhält man Energiekugeln an jedem Kontrollpunkt basierend auf der aktuellen Position.
	Auf der Strecke verteilt liegen Items, die beim hindurchfahren automatisch eingesammelt werden. Wer die Ziellinie zuerst überquert hat gewonnen. \par

	%Aussehen
	\codename\ hat einen futuristischen \gls{Low-Poly Stil}. Modelle besitzen keine Texturen, jedoch sollen die Materialien physikalisch akkurat dargestellt werden. Die Strecken
	verlaufen durch Raumstationen, welche hauptsächlich durch viele künstliche Lichtquellen beleuchtet werden. Die Stationen bestehen aus engen Tunneln und großen Kuppeln um Dynamik
	in das Spielgeschehen zu bringen. Die Checkpoints sollen für Spieler sichtbar und erkennbar dargestellt sein. Wichtige Momente während dem Rennen sollen durch visuelle Effekte hervorgehoben
	werden. \par

	\pagebreak
	
	\section{Zielbestimmung}

	%Grafik
	\subsection{InGame - Grafik}
	\subsubsection{Musskriterien}
	\begin{itemize}
		\item Darstellung physikalisch akkurater Materialien
		\item Rendern von undurchsichtigen Objekten
		\item Viele dynamische Lichtquellen gleichzeitig verarbeiten
		\item Schatten von dynamischen Lichtquellen
		\item Effekte
		\begin{itemize}
			\item Partikeleffekte (z.B. Explosionen)
		\end{itemize}
		\item Postprocessing (z.B. \gls{MotionBlur})
		\item Frei fliegende Kamera zum Ansehen der Strecke
	\end{itemize}
	
	\subsubsection{Wunschkriterien}
	\begin{itemize}
		\item Animationen an Spielfigur und Fahrzeug
		\begin{itemize}
			\item Klauen der Energie am vorderen Fahrzeug
			\item Animieren der Düsen des Fahrzeuges
		\end{itemize}
			\item Verschiedene Fahrzeuge
		\item Leuchtende Oberflächen
		\item \gls{HDR}
		\item \gls{FXAA}
		\item Soundeffekte
	\end{itemize}
	
	\pagebreak
	
	%Generirung
	\subsection{Prozedurale Generierung}
	\subsubsection{Musskriterien}
	\begin{itemize}
		\item Zusammenhängende Strecke mit Start und Ziel
		\item Checkpoints in sinnvollen Abständen
		\item Generierung von unterschiedlichen Streckenführungen
		\item Deterministisch vom \gls{Seed} abhängig
		\item Keine unbefahrbaren Streckenabschnitte
		\item Definieren von Extremparametern (Bsp. Steigung, Schärfe der Kurven, ...)
		\item Konform zum Streckenmodell des anderen Teams
		\item Generierung von Streckenumgebung
	\end{itemize}


	\subsubsection{Wunschkriterien}
	\begin{itemize}
		\item Verschiedene \glspl{Biom}
		\item Verschiedene befahrbare Untergründe
		\item Einfügen von Hindernissen auf die Rennbahn
		\item Dynamische Strecken- bzw. Umgebungselemente
	\end{itemize}
	
	\pagebreak
	
	%Menü, Pause, Hud
	\subsection{GUI}
	\subsubsection{Musskriterien} 
	\begin{itemize}
		\item Allgemeines Start-Menü
		\item Einstellungsfenster
		\begin{itemize}
			\item Diverse Möglichkeiten zur Konfiguration des Spielverhaltens
			\item Diverse Möglichkeiten zur Konfiguration der Grafik
		\end{itemize}
		\item Vor dem Spiel:
		\begin{itemize}
			\item Konfiguration des \glspl{Seed}
			\item Auswählen eines gespeicherten \glspl{Seed}
			\item Auswählen des Spielmodus
			\item Definition der Anzahl an optionaler \glspl{KI}
		\end{itemize}
		\item Ein während des Spiels aufrufbares Pause-Menü mit den Funktionen:
		\begin{itemize}
			\item Zurückkehren zum Spiel
			\item Zurückkehren zum Start-Menü
			\item Öffnen der Einstellungen
			\item Speichern des aktuellen \glspl{Seed}
		\end{itemize}
		\item Nach dem Spiel:
		\begin{itemize}
			\item Möglichkeit zur Speicherung des aktuellen \glspl{Seed}
			\item Eine Ergebnisanzeige mit Visualisierung der Endposition und der gefahrenen Zeit eines jeden Spielers
			\item Möglichkeit die Strecke erneut zu fahren
			\item Zurückkehren zum Start-Menü
		\end{itemize}
		\pagebreak
		\item Anzeige eines Head-up-Displays(HUD) während dem Spiel
		\begin{itemize}
			\item Mini-Map in 2D inklusive Visualisierung der Fahrer
			\item Anzeige der Rundenzahl
			\item Aktuell gefahrene Zeit
			\item Visualisierung des aufgesammelten Powerups
			\item Anzeige der Geschwindigkeit des Fahrers
			\item Anzeige der aktuellen Ranglistenplatzierung
		\end{itemize}
	\end{itemize}

	\subsubsection{Wunschkriterien}
	\begin{itemize}
		\item Konfigurationseinstellung für den Ton
		\item Fahrzeugauswahl
		\item Beispielvorschau eines \glspl{Seed}
		\item Abstimmmöglichkeit unter Spielern, ob eine Strecke erneut gespielt werden soll
	\end{itemize}

	\pagebreak
	
	\subsection{Assets}
	\subsubsection{Musskriterien}
	\begin{itemize}
		\item Verbindungstunnel der Raumstation
		\item Dekorationen für Raumstationselemente
		\item Kuppeln mit planetarischer Umgebung
		\item Hover-Cart (fliegendes Go-Cart)
		\item Energie-Kugeln hinten auf dem Hover-Cart
		\item Fahrer für Hover-Carts
	\end{itemize}
	
	\subsubsection{Wunschkriterien}
	\begin{itemize}
		\item Weitere Verbindungstunnel-Typen
		\item Weitere Fahrzeugmodelle
		\item Animation für Fahrer
	\end{itemize}
	
	\subsection{Abgrenzungskriterien}
	\begin{itemize}
		\item Das Spiel muss nicht auf Mobilgeräten oder Konsolen laufen
		
	\end{itemize}
	
	\pagebreak

	\section{Produkteinsatz}
	\subsection{Einsatzgebiet und Zielgruppe}
	Dieses Programm wird als Echtzeitcomputerspiel entwickelt.
	Dieses Spiel soll sich gezielt bei allen Familienmitgliedern etablieren.
	Lediglich eine Altersrestriktion von 6 Jahren aufwärts wird
	festgelegt. 
	Es soll gezielt der Gelegenheitsspieler angesprochen werden.
	Dabei sollte für die Vorkenntnisse des Benutzers die Altersrestriktion 
	beachtet werden. Es soll intuitiv benutzbar sein und einen einfachen 
	Zugang bieten. Lediglich minimalste Anforderungen an die Bedienung
	des Homecomputers sowie der Sprache werden vorausgesetzt.
	
	
	\subsection{Hard- und Softwareanforderungen}
	\begin{itemize}
		\item CPU: Intel i5 4.Generation aufwärts; AMD Ryzen 5 1600
		\item Grafikkarte: Nvidia GeForce GTX 960; Radeon RX 480 
		\item Eingabegeräte: Tastatur oder XBOX-Controller
		\item Arbeitsspeicher: 8GB
		\item Betriebssystem: Windows 10
	\end{itemize}
	\subsection{Laufzeit- und Speicherverbrauch}
	\begin{itemize}
		\item Rechenzeit pro Bild = 16ms (60bps)
		\item Speicherverbrauch maximal 3-4GB
		\item Interaktives Programm
	\end{itemize}
	
	\pagebreak
	
	\section{Anforderungen}

	\subsection{InGame - Grafik}
	\subsubsection{Funktionale Anforderungen}
	\begin{itemize}[label={}]
		\item \textbf{/FA000/} Darstellung physikalisch korrekter Materialien
		\item \textbf{/FA001/} Darstellung metallener bzw. dielektrischer Materialien
		\item \textbf{/FA002/} Darstellung untransparenter Objekte
		\item \textbf{/FA003/} Darstellung leuchtender Oberflächen
		
		\item \textbf{/FA010/} Darstellung von dynamischen Kegellichtquellen
		\item \textbf{/FA011/} Darstellung von statischen Kegellichtquellen
		\item \textbf{/FA012/} Darstellung von dynamischen Punktlichtquellen
		\item \textbf{/FA013/} Darstellung von statischen Punktlichtquellen
		\item \textbf{/FA014/} Darstellung von richtungsbasierten Lichtquellen(Sonne)
		\item \textbf{/FA015/} Berechnung  von dynamischen Schatten
		\item \textbf{/FA016/} Detailgrad der Schatten an Priorität der Lichter anpassen
		
		\item \textbf{/FA020/} Einbindung von Bildfiltern
		\item \textbf{/FA021/} Bewegungsunschärfe
		\item \textbf{/FA022/} Kantenglättungsfilter
		\item \textbf{/FA023/} Skalierung des Outputs der Renderpipeline
		\item \textbf{/FA024/} \gls{HDR}
		
		\item \textbf{/FA030/} Partikelsysteme
		\item \textbf{/FA031/} Variable Partikelgröße
		\item \textbf{/FA032/} Variable Partikelfarbe
		\item \textbf{/FA033/} Variable Partikelgeschwindigkeit
		\item \textbf{/FA034/} Beeinflussung der Partikel durch Gravitation
		\item \textbf{/FA035/} Funkenspritzen
		
		\item \textbf{/FA040/} Skelettbasierte Animation
		\item \textbf{/FA041/} Zweidimensionale Interpolation von Animationen
		\item \textbf{/FA042/} Animationsstatemachine
		\item \textbf{/FA043/} Schichten von Animationen
		\item \textbf{/FA044/} Animationsevents und Callbacks
		
		\item \textbf{/FA050/} Frei fliegende Kamera
	\end{itemize}

	
	\subsubsection{Nichtfunktionale Anforderungen}
	\begin{itemize}[label={}]
		\item \textbf{/NFA000/} Konstante Bildrate (ca. 60bps)
		\item \textbf{/NFA010/} Mindestens 4 gleichzeitig aktive Lichter
		\item \textbf{/NFA011/} Berechnung von dynamischen Schatten für mindestens 2 Lichter gleichzeitig
		\item \textbf{/NFA020/} Mindestens 2 aktive Bildfilter gleichzeitig
		\item \textbf{/NFA030/} Darstellung von mindestens 100 Partikeln
		\item \textbf{/NFA040/} Skelette mit mindestens 26 Knochen
		\item \textbf{/NFA041/} Mindestens 2 Schichten von Animationen
	\end{itemize}
	
	\pagebreak
	
	\subsection{Prozedurale Generierung}
	\subsubsection{Funktionale Anforderungen}
	\begin{itemize} [label={}]
		\item \textbf{/FA100/} Generierung einer durchgehenden Strecke mit Start und Ziel
		\item \textbf{/FA101/} Ein \gls{Seed} produziert deterministisch immer die gleiche Ausgabe des Generators
		\item \textbf{/FA102/} Ausgabe von \gls{Navigation Mesh}
		\item \textbf{/FA103/} Ausgabe von \gls{Scenegraph}
		
		\item \textbf{/FA110/} Festlegen von Checkpoints
		\item \textbf{/FA111/} Visuelle Repräsentation von Start, Ziel und Checkpoints
		\item \textbf{/FA112/} Generieren der Startpositionen für Fahrzeuge
		\item \textbf{/FA113/} Generieren der Positionen für Powerups
		\item \textbf{/FA114/} Generieren von unterschiedlichen Fahruntergründen
		
		\item \textbf{/FA120/} Unterteilung der Strecke in \glspl{Biom}
		\item \textbf{/FA121/} Generieren der Streckenumgebung
		\item \textbf{/FA122/} Generieren von Objekten auf der Strecke
		
		\item \textbf{/FA130/} Algorithmus durch verschiedene Parameter konfigurierbar
		\item \textbf{/FA131/} Berücksichtigung von Extremparametern
	\end{itemize}
	
	\subsubsection{Nichtfunktionale Anforderungen}
	\begin{itemize} [label={}]
		\item \textbf{/NFA100/} Checkpoints in gleichmäßigen Abständen
		
		\item \textbf{/NFA110/} Abwechslungsreiches Fahrerlebnis innerhalb einer Strecke
		\item \textbf{/NFA111/} Strecken sollen sich für den Spieler merklich Unterscheiden, bei unterschiedlichen \glspl{Seed}
		
		\item \textbf{/NFA120/} Strecke muss durchgehend bis zum Ziel befahrbar sein
		\item \textbf{/NFA121/} Die Generierung muss in unter 30 Sekunden erfolgen
	\end{itemize}
		
	\pagebreak
	
	\subsection{GUI}
	\subsubsection{Funktionale Anforderungen}
	\begin{itemize} [label={}]
		\item \textbf{/FA200/} Das Spiel startet mit einem auf dem Bildschirm erscheinenden Start-Menü
		\item \textbf{/FA201/} Durch Drücken eines Buttons lässt sich das Einstellungsfenster öffnen
		\item \textbf{/FA203/} Das Spiel kann im Einzelspielermodus gestartet werden
		\item \textbf{/FA204/} Das Spiel kann im Mehrspielermodus gestartet werden
		
		\item \textbf{/FA210/} Im Einstellungsfenster können spielinterne Einstellungen vorgenommen werden
		\item \textbf{/FA211/} Im Einstellungsfenster können grafische Einstellungen vorgenommen werden
		\item \textbf{/FA212/} Im Einstellungsfenster kann der Spieler zwischen Vollbild- und Fenstermodus wechseln
		\item \textbf{/FA213/} In den spielinternen Einstellungen kann die Steuerung eingestellt werden
		
		\item \textbf{/FA220/} Nach Auswählen des Einzelspielermodus öffnet sich ein neues Fenster zur Streckenauswahl
		\item \textbf{/FA221/} In der Streckenauswahl kann der \gls{Seed} konfiguriert werden werden	
		\item \textbf{/FA222/} In der Streckenauswahl kann ein gespeicherter \gls{Seed} ausgewählt werden
		\item \textbf{/FA223/} In der Streckenauswahl kann der Name eines gespeicherten \glspl{Seed} bearbeitet werden
		\item \textbf{/FA224/} In der Streckenauswahl kann im Einzelspielermodus die Anzahl an \glspl{KI} eingestellt werden	
		\item \textbf{/FA225/} In der Streckenauswahl kann das Spiel durch Drücken eines Buttons gestartet werden
		
		\item \textbf{/FA230/} Nach Auswählen des Mehrspielermodus öffnet sich eine Lobby
		\item \textbf{/FA231/} In der Lobby werden alle bisher verbundenen Spieler angezeigt
		\item \textbf{/FA232/} In der Lobby existiert eine Fahrzeugauswahl
		\item \textbf{/FA233/} In der Lobby kann man ein Einstellungsfenster öffnen, in dem man spielinterne Einstellungen vornehmen kann
		\item \textbf{/FA234/} In der Lobby kann das Spiel gestartet werden
		\item \textbf{/FA235/} In der Lobby besteht die Möglichkeit den Namen der Lobby zu ändern
		\item \textbf{/FA236/} Eine Lobby kann per Knopfdruck verlassen wiederverwendbar
		\item \textbf{/FA237/} Der Spieler kann einer anderen Lobby beitreten
		
		\item \textbf{/FA240/} Im Spiel kann durch Drücken des entsprechenden Knopfes das Pause-Menü geöffnet werden
		\item \textbf{/FA241/} Im Pause-Menü kann man zum Spiel zurückkehren
		\item \textbf{/FA242/} Im Pause-Menü kann man zum Start-Menü wechseln
		\item \textbf{/FA243/} Im Pause-Menü kann man das Einstellungsfenster öffnen
		\item \textbf{/FA244/} Im Pause-Menü kann man den aktuellen \glspl{Seed} speichern
		
		\item \textbf{/FA250/} Im Spiel wird die Geschwindigkeit des eigenen Fahrzeugs an diesem angezeigt
		\item \textbf{/FA251/} Die Anzahl an Energiekugeln ist am Heck des Fahrzeuges zu sehen
		\item \textbf{/FA252/} Die Position und Fahrzeit wird angezeigt
		\item \textbf{/FA253/} Eine Mini-Map mit den Positionen aller Fahrer wird dargestellt
		\item \textbf{/FA254/} Das momentan ausgerüstete Power-Up wird dargestellt
		
		\item \textbf{/FA260/} Erreicht der Spieler das Ziel, erscheint eine Ergebnisanzeige, die eine Rangliste aller Fahrer anzeigt sowie deren Fahrzeiten
		\item \textbf{/FA261/} In der Ergebnisanzeige kann der \glspl{Seed} gespeichert werden
		\item \textbf{/FA262/} In der Ergebnisanzeige kann durch Drücken eines Buttons die Strecke erneut gefahren werden
		\item \textbf{/FA263/} In der Ergebnisanzeige besteht die Möglichkeit zum Start-Menü zurückzukehren
	\end{itemize}

	\subsubsection{Nichtfunktionale Anforderungen}
	\begin{itemize} [label={}]
		\item \textbf{/NFA200/} Das Start-Menü öffnet sich innerhalb von 30 Sekunden
		\item \textbf{/NFA201/} Jegliche Menüs reagieren nur auf gültige Eingaben/Klicks
		\item \textbf{/NFA202/} Jegliche Benutzereinstellungen werden korrekt angewendet
		\item \textbf{/NFA203/} Der End-Bildschirm sollte sich nach Ende des Rennens öffnen
	\end{itemize}

	\pagebreak
		
	\subsection{Assets}
	\subsubsection{Funktionale Anforderungen}
	\begin{itemize} [label={}]
		\item \textbf{/FA300/} Modelle für Tunnel-Segmente
		\item \textbf{/FA301/} Modelle für Dekorationselmente
		\item \textbf{/FA302/} Modelle für Kuppel-Ränder
		\item \textbf{/FA303/} Modelle für Vegetation innerhalb der Kuppeln
		\item \textbf{/FA310/} Modelle für Hover-Carts
		\item \textbf{/FA311/} Modelle für Energie-Kugeln
		\item \textbf{/FA312/} Modelle für Fahrer
		\item \textbf{/FA313/} Animationen für Fahrer
		\item \textbf{/FA314/} Animationen für Hover-Carts
	\end{itemize}
	
	\pagebreak
	
	\section{Systemmodell}
	
	\begin{figure}[!h]
	\centering
    \centering
    	\includegraphics[width=.85\linewidth]{./Bilder/Systemmodell_Hauptmenu.png}
  		\caption{Systemmodell des Hauptmenüs}
  	 	\label{fig:Systemmodell_Hauptmenu}
  	 	
  	 	\vspace{0.7cm}
  	 	
    	\includegraphics[width=.85\linewidth]{./Bilder/Systemmodell_Multiplayer_Lobby.png}
  		\caption{Systemmodell der Multiplayer-Lobby}
  	 	\label{fig:Systemmodell_Multiplayer_Lobby}

	\end{figure}
	
	\begin{figure}[h]
	\centering
	\centering
    	\includegraphics[width=.85\linewidth]{./Bilder/Systemmodell_System_Kommunikation.png}
  		\caption{Systemmodell vom Systemkommunikation}
  	 	\label{fig:Systemmodell_System_Kommunikation}
	\end{figure}	 
	 
	%wegen Einfügen der Bilder so nötig
	\newpage
	\ \\
	\newpage
	 
	\section{Bedienoberfläche}
	
	\begin{center}
    	\includegraphics[width=.58\linewidth]{./Bilder/startmenue.png}
  		\captionof{figure}{Startmenü}
  	 	\label{fig:Startmenue}		
	
    	\includegraphics[width=.58\linewidth]{./Bilder/streckenauswahl_einzel.png}
  		\captionof{figure}{Streckenauswahl im Einzelspielermodus}
  	 	\label{fig:Streckenauswahl_Einzel}

    	\includegraphics[width=\linewidth]{./Bilder/lobby.png}
  		\captionof{figure}{Lobby im Mehrspielermodus}
  	 	\label{fig:Lobby}

    	\includegraphics[width=\linewidth]{./Bilder/spielinterne_einstellungen.png}
  		\captionof{figure}{Spielinterne Einstellungen in der Lobby}
		   \label{fig:Spielinterne_einstellungen}

		\includegraphics[width=\linewidth]{./Bilder/HudMockUp.png}
		\captionof{figure}{HUD während des Rennens}
			\label{fig:HUD}

    	\includegraphics[width=.4\linewidth]{./Bilder/pausemenue.png}
  		\captionof{figure}{Pausemenü}
  	 	\label{fig:Pausemenue}
  	 	
  	 	\includegraphics[width=.6\linewidth]{./Bilder/einstellungen.png}
  		\captionof{figure}{Einstellungsfenster}
  	 	\label{fig:Einstellungen}

    	\includegraphics[width=.8\linewidth]{./Bilder/endscreen.png}
  		\captionof{figure}{Endbildschirm}
  	 	\label{fig:Endbildschirm}
	\end{center}
	
	\newpage
	
	\section{Qualitätssicherung}
	
	Für eine möglichst umfassende Qualitätssicherung wird der gesamte Testablauf in zwei verschiedene Arten aufgeteilt.
	Die \textit{internen Testabläufe}, welche Fehler aufdecken sollen, die direkt im Programmcode auftreten.
	Des weiteren dienen die \textit{externen Testabläufe} dazu, unbeabsichtigtes Verhalten der Anwendung aufzudecken wie beispielsweise inkonsistente Menüführung
	und für das Rennen ungeeignete Parameter in der Generierung.
	
	\subsection{Interner Testablauf}
	
	Durch die \textit{internen Testabläufe} soll sichergestellt werden, dass das Projekt im Programmcode möglichst wenige Fehler enthält und außerdem möglichst sicher ist,
	um beispielsweise Probleme bei falschen Konfigurationen zu vermeiden.\par
	
	Um dies zu erreichen wird unser Code schon während der Entwicklung getestet, mit Werkzeugen wie beispielsweise \gls{JUnit}. Hierbei wird eine Testabdeckung von mindestens 80\% angestrebt.
	Ausgenommen ist jedoch Programmcode, welcher auf der Grafikkarte ausgeführt wird, da dieser sich nicht mit \gls{JUnit} testen lässt.
	Die Tests beinhalten dann beispielsweise falsche Konfigurationen und ungewöhnliche Eingaben, um möglichst alle Randfälle des Programms zu überprüfen.
	
	
	\subsection{Externe Testabläufe}
	
	Die \textit{externen Testabläufe} werden verwendet um Fehler zu finden, welche nicht durch automatisierte Tests erkannt werden können. Um dies zu erreichen wird das Spiel in seiner Ausführung beobachtet
	und analysiert.
	Um möglichst viele dieser Fehler beheben zu können werden diese Tests in wohl überlegten Szenarien und Abläufen durchgeführt, welche im Folgenden definiert werden.\par
	
	Nachdem diese Tests alle absolviert wurden wird das Spiel des weiteren ohne konkrete Vorgaben gespielt, also Tests von Anwendungsstart bis zum Beenden des Spiels durchgeführt.
	Da dies ohne Vorgaben geschieht, versucht man so anschließend Fehlerquellen zu finden, an die zuvor möglicherweise nicht gedacht wurde.
	
	\pagebreak
	
	\subsubsection{InGame - Grafik}
	\begin{itemize} [label={}]
		\item \textbf{ET010}\par
			\textbf{STAND} Szene mit dynamischem Licht.\\
			\textbf{AKTION} Licht wird bewegt.\\
			\textbf{RESULTAT} Beleuchtung der Szene ist der neuen Lichtposition angepasst.\\\par
		
		\item \textbf{ET011}\par
			\textbf{STAND} Szene mit dynamischem Licht.\\ 
			\textbf{AKTION} Objekt wird bewegt.\\
			\textbf{RESULTAT} Beleuchtung der Szene ist der neuen Objektposition angepasst.\\\par

		\item \textbf{ET012}\par
			\textbf{STAND} Szene mit dynamischem Licht.\\
			\textbf{AKTION} Kamera wird bewegt.\\
			\textbf{RESULTAT} Perspektive wird der neuen Kameraposition angepasst. Lichteffekte (wie Reflektionen) werden ebenfalls angepasst.\\\par
		
		\item \textbf{ET013}\par
			\textbf{STAND} Szene mit dynamischen Lichtern.\\
			\textbf{AKTION} Kamera wird bewegt.\\
			\textbf{RESULTAT} Menge der Lichter, für welche Schatten berechnet werden, wird angepasst.\\\par

		\item \textbf{ET020}\par
			\textbf{STAND} Szene mit dynamischen Lichtern.\\
			\textbf{AKTION} Bildeffekt wird hinzugefügt.\\
			\textbf{RESULTAT} Bildeffekt wird angewendet.\\\par

		\item \textbf{ET030}\par
			\textbf{STAND} Szene mit inaktivem Partikel-Emitter.\\
			\textbf{AKTION} Partikel-Emitter wird aktiviert.\\
			\textbf{RESULTAT} Animiertes Partikelsystem ist zu sehen.\\\par

		\pagebreak
		\item \textbf{ET040}\par
			\textbf{STAND} Objekt, welches animiert werden kann.\\
			\textbf{AKTION} Animation wird gestartet.\\
			\textbf{RESULTAT} Objekt ist animiert und Animationsevents werden aufgerufen.\\\par
	\end{itemize}

	\subsubsection{Prozedurale Generierung}
	\begin{itemize} [label={}]
		\item \textbf{ET100}\par
			\textbf{STAND} Startmenü.\\ 
			\textbf{AKTION} Benutzer startet Strecken-Generierung mit beliebigem \gls{Seed}.\\
			\textbf{RESULTAT} Strecke wird generiert.\\\par
		
		\item \textbf{ET101}\par
			\textbf{STAND} Startmenü.\\ 
			\textbf{AKTION} Benutzer startet Strecken-Generierung nacheinander min. 2 mal mit dem gleichen \gls{Seed} und analysiert anschließend die Strecke.\\
			\textbf{RESULTAT} Die generierten Strecken müssen alle identisch sein.\\\par
			
		\item \textbf{ET102}\par
			\textbf{STAND} Startmenü.\\ 
			\textbf{AKTION} Benutzer startet Strecken-Generierung nacheinander min. 2 mal mit paarweise verschiedenen \glspl{Seed} und analysiert anschließend die Strecke.\\
			\textbf{RESULTAT} Alle generierten Strecken müssen sich unterscheiden.\\\par
		
		\item \textbf{ET110}\par
			\textbf{Stand} Strecke wurde fertig generiert.\\
			\textbf{AKTION} Strecke wird überquert und der grobe Aufbau analysiert.\\
			\textbf{RESULTAT} Strecke ist zusammenhängend (mit Start und Ziel).\\\par
			
		\item \textbf{ET111}\par
			\textbf{Stand} Strecke wurde fertig generiert.\\
			\textbf{AKTION} Strecke wird überquert und der grobe Aufbau analysiert.\\
			\textbf{RESULTAT} Alle Checkpoints stehen auf der Strecke und in regelmäßigen Abständen.\\\par
			
		\pagebreak
		\item \textbf{ET112}\par
			\textbf{Stand} Strecke wurde fertig generiert.\\
			\textbf{AKTION} Strecke wird überquert und die Fahrbahn analysiert.\\
			\textbf{RESULTAT} Fahrbahn muss gut befahrbar sein, d.h. innerhalb der Extremparameter sein.\\\par
			
		\item \textbf{ET113}\par
			\textbf{STAND} Strecke wurde fertig generiert.\\
			\textbf{AKTION} Strecke wird überquert und in Bezug auf die Umgebung analysiert.\\
			\textbf{RESULTAT} Strecke interagiert sinnvoll mit der Umgebung (d.h. Strecke schwebt nicht ungewollt über oder unter dem Boden).\\\par
			
		\item \textbf{ET114}\par
			\textbf{STAND} Strecke wurde fertig generiert.\\ 
			\textbf{AKTION} Strecke wird überquert und Objekte auf und um die Fahrbahn beobachtet.\\
			\textbf{RESULTAT} Es dürfen keine Objekte in die Fahrbahn ragen, welche dort nicht sein sollen. Außerdem dürfen Objekte auf der Fahrbahn diese nicht gänzlich blockieren (unbefahrbar machen).\\\par
			
		
	\end{itemize}

	\subsubsection{GUI}
		\begin{itemize} [label={}]
			
		\item \textbf{ET200}\par
			\textbf{STAND} Startmenü\\ 
			\textbf{AKTION} Benutzer navigiert zum Einzelspielermodus, wählt eine gespeicherte Strecke oder generiert eine Neue und startet das Spiel.\\
			\textbf{RESULTAT} Spiel wird mit entsprechender Strecke gestartet.\\\par	
		
		\item \textbf{ET201}\par
			\textbf{STAND} Startmenü\\ 
			\textbf{AKTION} Der Benutzer navigiert ins Einstellungsmenü, nimmt einige Einstellungen vor und überprüft diese im Spiel.\\
			\textbf{RESULTAT} Einstellungen wurden gespeichert und werden angewendet.\\\par
		
		\item \textbf{ET202}\par
			\textbf{STAND} Startmenü\\ 
			\textbf{AKTION} Benutzer startet eine Multiplayerlobby mit mindestens einem Mitspieler.\\
			\textbf{RESULTAT} Lobbymenü wird aufgerufen und zeigt alle Einstellungen und Spieler korrekt an.\\\par
		
		\item \textbf{ET210}\par
			\textbf{STAND} Einstellungsmenü\\ 
			\textbf{AKTION} Der Benutzer nimmt Einstellungen im Einstellungsmenü vor, startet das Spiel anschließend neu und überprüft, ob Einstellungen angewandt werden.\\
			\textbf{RESULTAT} Einstellungen wurden gespeichert und werden angewendet.\\\par
			
		\item \textbf{ET220}\par
			\textbf{STAND} Beliebiges Menü\\
			\textbf{AKTION} Der Benutzer navigiert über die entsprechenden Knöpfe ins Hauptmenü.\\
			\textbf{RESULTAT} Hauptmenü wird aufgerufen und ggf. das laufende Rennen verlassen.\\\par
		
		\item \textbf{ET230}\par
			\textbf{STAND} Im laufenden Rennen\\ 
			\textbf{AKTION} Der Benutzer ruft das Pausemenü auf und verwendet dieses.\\
			\textbf{RESULTAT} Das Menü angezeigt. Entsprechende Aktionen werden ausgeführt.\\\par
			
		\item \textbf{ET240}\par
			\textbf{STAND} Im laufenden Rennen\\ 
			\textbf{AKTION} Der Benutzer fährt das Rennen und beobachtet die Anzeigen des HUD.\\
			\textbf{RESULTAT} Alle Anzeigen entsprechen dem Spielgeschehen.\\\par
	
		\end{itemize}

	
	\pagebreak
	
	\section{Glossar}
	\printnoidxglossaries
	
\end{document}
